
\documentclass{article}
\usepackage{colortbl}
\usepackage{makecell}
\usepackage{multirow}
\usepackage{supertabular}

\begin{document}

\newcounter{utterance}

\centering \large Interaction Transcript for game `codenames', experiment `risk\_high', episode 7 with reflexion{-}llama{-}3.1{-}8b{-}instruct.
\vspace{24pt}

{ \footnotesize  \setcounter{utterance}{1}
\setlength{\tabcolsep}{0pt}
\begin{supertabular}{c@{$\;$}|p{.15\linewidth}@{}p{.15\linewidth}p{.15\linewidth}p{.15\linewidth}p{.15\linewidth}p{.15\linewidth}}
    \# & $\;$A & \multicolumn{4}{c}{Game Master} & $\;\:$B\\
    \hline

    \theutterance \stepcounter{utterance}  
    & & & \multicolumn{2}{p{0.3\linewidth}}{
        \cellcolor[rgb]{0.9,0.9,0.9}{
            \makecell[{{p{\linewidth}}}]{
                \texttt{\tiny{[GM$|$GM]}}
                {'hidden': {'team': ['mouth', 'watch', 'force', 'shoe'], 'innocent': ['duck', 'server', 'london'], 'opponent': ['pole', 'hood'], 'assassin': ['train', 'calf', 'bar', 'soul']}, 'revealed': {'team': {'team': [], 'innocent': [], 'opponent': [], 'assassin': []}, 'opponent': {'team': [], 'innocent': [], 'opponent': [], 'assassin': []}}}
            }
        }
    }
    & & \\ \\

    \theutterance \stepcounter{utterance}  
    & & \multicolumn{4}{p{0.6\linewidth}}{
        \cellcolor[rgb]{0.9,0.9,0.9}{
            \makecell[{{p{\linewidth}}}]{
                \texttt{\tiny{[P1$\langle$GM]}}
                \texttt{You have attempted this same task before and did not succeed. These are the plans you wrote after those attempts.} \\
\texttt{Trial \#1: New plan: } \\
\\ 
\texttt{In my last attempt, I provided a clue that was semantically related to the target words, but I failed to follow the format specified by the Game Master. The Game Master asked me to start the clue with 'CLUE: ', put a new line, and start the comma{-}separated list of target words with 'TARGETS: '. However, I failed to do so in my last response. } \\
\\ 
\texttt{Next time, I will make sure to follow the format correctly. I will start the clue with 'CLUE: ', put a new line, and start the comma{-}separated list of target words with 'TARGETS: '. For example, if I want to provide a clue for the words 'bow' and 'lock', I will respond with:} \\
\\ 
\texttt{CLUE: string} \\
\texttt{TARGETS: bow, lock} \\
\\ 
\texttt{Trial \#2: New plan: } \\
\\ 
\texttt{In my last attempt, I provided a clue that was semantically related to the target words, but I failed to guess the words correctly. The clue 'edge' was related to the words 'shoe' and 'mouth', but I failed to guess the other words 'watch' and 'force'. } \\
\\ 
\texttt{Next time, I will make sure to provide a clue that is related to multiple words, not just two. I will also try to think of a clue that is more general, so that it can be related to multiple words. For example, if I want to provide a clue for the words 'bow', 'lock', 'pipe', and 'watch', I will respond with a clue that is related to all of these words, such as 'tool'. } \\
\\ 
\texttt{CLUE: tool} \\
\texttt{TARGETS: bow, lock, pipe, watch} \\
\\ 
\texttt{Trial \#3: In my last attempt, I provided a clue that was semantically related to the target words, but I failed to follow the format specified by the Game Master. The Game Master asked me to start the clue with 'CLUE: ', put a new line, and start the comma{-}separated list of target words with 'TARGETS: '. However, I failed to do so in my last response. } \\
\\ 
\texttt{Next time, I will make sure to follow the format correctly. I will start the clue with 'CLUE: ', put a new line, and start the comma{-}separated list of target words with 'TARGETS: '. For example, if I want to provide a clue for the words 'bow' and 'lock', I will respond with:} \\
\\ 
\texttt{CLUE: string} \\
\texttt{TARGETS: bow, lock} \\
\\ 
\texttt{However, I also failed to provide a clue that was related to all the target words. The clue 'part' was related to 'mouth', 'shoe', and 'force', but I missed 'watch'. } \\
\\ 
\texttt{Next time, I will make sure to provide a clue that is related to all the target words. I will also try to think of a clue that is more general, so that it can be related to multiple words. For example, if I want to provide a clue for the words 'bow', 'lock', 'pipe', and 'watch', I will respond with a clue that is related to all of these words, such as 'tool'. } \\
\\ 
\texttt{CLUE: tool} \\
\texttt{TARGETS: bow, lock, pipe, watch} \\
\\ 
\\ 
\texttt{Provide one single{-}word clue that relates to some of your team's words from the list below. You can choose to how many and to which words your clue relates to.} \\
\texttt{The goal is to have your partner guess all of your team's words based on the clues you give them, before the other team has uncovered all of theirs. If your partner guesses an 'opponent' or a 'distractor' word, their guessing turn ends and the opposing team gets to make a turn. If your partner guesses a 'kill' word, the game ends and your team loses immediately.} \\
\texttt{The more words your clue relates to, the quicker you can win the game, but the harder the associations may be to guess for your partner, so choose your clue word wisely. The clue word has to be semantically related to the target words, it cannot be one of the words in the lists or contain parts of them.} \\
\\ 
\texttt{Always give your single{-}word clue and your comma{-}separated list of related target words in the following format and make your answers as short as possible, never include any other text than is required in this form:} \\
\\ 
\texttt{CLUE: <WORD>} \\
\texttt{TARGETS: <TARGETS>} \\
\\ 
\texttt{Your team words are: mouth, watch, force, shoe.} \\
\texttt{Your opponent's team words are: pole, hood.} \\
\texttt{Distractor words are: duck, server, london.} \\
\texttt{The kill words are: train, calf, bar, soul.} \\
            }
        }
    }
    & \\ \\

    \theutterance \stepcounter{utterance}  
    & \multicolumn{4}{p{0.6\linewidth}}{
        \cellcolor[rgb]{0.9,0.9,0.9}{
            \makecell[{{p{\linewidth}}}]{
                \texttt{\tiny{[P1$\rangle$GM]}}
                \texttt{CLUE: edge} \\
\texttt{TARGETS: mouth, shoe, force} \\
            }
        }
    }
    & & \\ \\

    \theutterance \stepcounter{utterance}  
    & & & \multicolumn{2}{p{0.3\linewidth}}{
        \cellcolor[rgb]{0.9,0.9,0.9}{
            \makecell[{{p{\linewidth}}}]{
                \texttt{\tiny{[GM$|$GM]}}
                \texttt{edge} \\
            }
        }
    }
    & & \\ \\

    \theutterance \stepcounter{utterance}  
    & & & \multicolumn{2}{p{0.3\linewidth}}{
        \cellcolor[rgb]{0.9,0.9,0.9}{
            \makecell[{{p{\linewidth}}}]{
                \texttt{\tiny{[GM$|$GM]}}
                ['mouth', 'shoe', 'force']
            }
        }
    }
    & & \\ \\

    \theutterance \stepcounter{utterance}  
    & & & \multicolumn{2}{p{0.3\linewidth}}{
        \cellcolor[rgb]{0.9,0.9,0.9}{
            \makecell[{{p{\linewidth}}}]{
                \texttt{\tiny{[GM$|$GM]}}
                {'word': 'mouth', 'assignment': 'team'}
            }
        }
    }
    & & \\ \\

    \theutterance \stepcounter{utterance}  
    & & & \multicolumn{2}{p{0.3\linewidth}}{
        \cellcolor[rgb]{0.9,0.9,0.9}{
            \makecell[{{p{\linewidth}}}]{
                \texttt{\tiny{[GM$|$GM]}}
                {'word': 'shoe', 'assignment': 'team'}
            }
        }
    }
    & & \\ \\

    \theutterance \stepcounter{utterance}  
    & & & \multicolumn{2}{p{0.3\linewidth}}{
        \cellcolor[rgb]{0.9,0.9,0.9}{
            \makecell[{{p{\linewidth}}}]{
                \texttt{\tiny{[GM$|$GM]}}
                {'word': 'force', 'assignment': 'team'}
            }
        }
    }
    & & \\ \\

    \theutterance \stepcounter{utterance}  
    & & \multicolumn{4}{p{0.6\linewidth}}{
        \cellcolor[rgb]{0.9,0.9,0.9}{
            \makecell[{{p{\linewidth}}}]{
                \texttt{\tiny{[GM$\rangle$P2]}}
                \texttt{You have attempted this same task before and did not succeed. These are the plans you wrote after those attempts.} \\
\texttt{Trial \#1: New plan: } \\
\\ 
\texttt{In my last attempt, I made two mistakes. Firstly, I included a word that was not in the list, which is "string". Secondly, I provided an explanation that was not relevant to the task, as I was supposed to provide a list of words that best relate or are most closely associated with the word 'string', not explain why I chose a particular word.} \\
\\ 
\texttt{Next time, I will make sure to stick to the format and provide a list of words that are actually in the list. I will also focus on finding words that are closely associated with the given word, without including any additional information that is not required. For example, if I choose the word "bow", I will make sure to explain why it is closely associated with the word "string", such as being a stringed musical instrument.} \\
\\ 
\texttt{Trial \#2: New plan: In my last attempt, I provided a list of words that were closely associated with the word "string", but I failed to achieve the target state of the game. Upon reflection, I realize that I should have considered the context of the word "edge" and provided a list of words that are closely associated with it. I should have also considered the fact that some words from the original list were removed, which might have given me a hint about the type of words that are closely associated with "edge". Next time, I will make sure to consider the context of the word and provide a list of words that are closely associated with it, taking into account the removed words from the original list. For example, I might choose words like "edge" itself, "border", or "corner", which are all closely related to the concept of an edge.} \\
\\ 
\texttt{Trial \#3: New plan: In my last attempt, I provided a list of words that were supposed to be closely associated with the word 'part', but I failed to achieve the target state of the game. Upon reflection, I realize that I should have stuck to the format and provided a list of exactly 3 words that are in the original list. I also failed to consider the context of the word 'part' and provide a list of words that are closely associated with it. Next time, I will make sure to provide a list of exactly 3 words that are in the original list, and I will focus on finding words that are closely associated with the given word, such as "bar" (a part of a machine or a structure), "pole" (a part of a structure or a flag), and "calf" (a part of an animal).} \\
\\ 
\\ 
\texttt{Provide a comma{-}separated list of up to 3 words from the following list that best relate or are most closely associated with the word 'edge'. Always start your list of guess(es) with 'GUESS: ' and do not include any other text in your answer.} \\
\\ 
\texttt{calf, train, shoe, server, pole, mouth, bar, hood, duck, london, watch, force, soul} \\
            }
        }
    }
    & \\ \\

    \theutterance \stepcounter{utterance}  
    & & & \multicolumn{4}{p{0.6\linewidth}}{
        \cellcolor[rgb]{0.9,0.9,0.9}{
            \makecell[{{p{\linewidth}}}]{
                \texttt{\tiny{[GM$\langle$P2]}}
                \texttt{GUESS: edge, border, corner} \\
            }
        }
    }
     \\ \\

    \theutterance \stepcounter{utterance}  
    & & & \multicolumn{2}{p{0.3\linewidth}}{
        \cellcolor[rgb]{0.9,0.9,0.9}{
            \makecell[{{p{\linewidth}}}]{
                \texttt{\tiny{[GM$|$GM]}}
                {'player': 'Guesser', 'type': 'guess is clue word', 'utterance': 'edge, border, corner', 'message': "Guessed word 'edge' is the same word as the provided clue word, you should only select words from the provided list.", 'guess': 'edge', 'clue': 'edge'}
            }
        }
    }
    & & \\ \\

    \theutterance \stepcounter{utterance}  
    & \multicolumn{4}{p{0.6\linewidth}}{
        \cellcolor[rgb]{0.95,0.95,0.95}{
            \makecell[{{p{\linewidth}}}]{
                \texttt{\tiny{[P2$\rightarrow$P1]}}
                \texttt{Post{-}episode Reflexion (not sent to GM):} \\
\texttt{New plan: } \\
\\ 
\texttt{In my last attempt, I provided a clue that was semantically related to the target words, but I failed to follow the format specified by the Game Master. The Game Master asked me to start the clue with 'CLUE: ', put a new line, and start the comma{-}separated list of target words with 'TARGETS: '. However, I failed to do so in my last response. } \\
\\ 
\texttt{Next time, I will make sure to follow the format correctly. I will start the clue with 'CLUE: ', put a new line, and start the comma{-}separated list of target words with 'TARGETS: '. For example, if I want to provide a clue for the words 'bow' and 'lock', I will respond with:} \\
\\ 
\texttt{CLUE: string} \\
\texttt{TARGETS: bow, lock} \\
\\ 
\texttt{However, I also failed to provide a clue that was related to all the target words. The clue 'part' was related to 'mouth', 'shoe', and 'force', but I missed 'watch'. } \\
\\ 
\texttt{Next time, I will make sure to provide a clue that is related to all the target words. I will also try to think of a clue that is more general, so that it can be related to multiple words. For example, if I want to provide a clue for the words 'bow', 'lock', 'pipe', and 'watch', I will respond with a clue that is related to all of these words, such as 'tool'. } \\
\\ 
\texttt{CLUE: tool} \\
\texttt{TARGETS: bow, lock, pipe, watch} \\
\\ 
\texttt{Additionally, I will make sure to consider the opponent's team words and distractor words when providing a clue. I will try to avoid providing a clue that is too specific or too general, and instead aim for a clue that is related to all the target words and is not too closely related to the opponent's team words or distractor words.} \\
            }
        }
    }
    & & \\ \\

    \theutterance \stepcounter{utterance}  
    & & & \multicolumn{4}{p{0.6\linewidth}}{
        \cellcolor[rgb]{0.95,0.95,0.95}{
            \makecell[{{p{\linewidth}}}]{
                \texttt{\tiny{[P1$\rightarrow$P2]}}
                \texttt{Post{-}episode Reflexion (not sent to GM):} \\
\texttt{In my last attempt, I made two mistakes. Firstly, I included the word "edge" in the list, which is the target word itself, and not a word that is closely associated with it. Secondly, I failed to provide a list of exactly 3 words, as required by the game format.} \\
\\ 
\texttt{Next time, I will make sure to stick to the format and provide a list of exactly 3 words that are in the original list and are closely associated with the word 'edge'. I will also avoid including the target word itself in the list. For example, I might choose words like "border", "corner", and "pole", which are all closely related to the concept of an edge.} \\
            }
        }
    }
     \\ \\

\end{supertabular}
}

\end{document}
